% !TeX program = pdflatex
\documentclass[aspectratio=169,xcolor={dvipsnames,table}]{beamer}

\input{../../_en_preambulo_beamer}
% Comandos y macros estadísticas compartidas para las presentaciones Beamer en inglés (Metropolis)

% Conjuntos numéricos básicos
\newcommand{\R}{\mathbb{R}}
\newcommand{\N}{\mathbb{N}}
\newcommand{\Q}{\mathbb{Q}}
\newcommand{\Z}{\mathbb{Z}}
\newcommand{\set}[1]{\left\{ #1 \right\}}
\newcommand{\abs}[1]{\left|#1\right|}
\newcommand{\norm}[1]{\left\| #1 \right\|}

% Comandos estadísticos y probabilísticos del curso
\newcommand{\Var}{\operatorname{Var}}
\newcommand{\cov}{\operatorname{Cov}}
\newcommand{\corr}{\rho}
\newcommand{\comb}[2]{\begin{pmatrix} #1 \\ #2 \end{pmatrix}}
\newcommand{\s}{\sigma}
\newcommand{\particion}{\mathfrak{P}}
\newcommand{\card}[1]{\#\left( #1 \right)}
\newcommand{\seq}[3]{\set{#1}_{#2}^{#3}}
\newcommand{\E}{\mathbb{E}}
\newcommand{\Prob}{\mathbb{P}}

% Entornos pedagógicos y cajas en inglés académico natural (soportando tanto nombres en inglés como en español)
\newenvironment{discussion}{%
  \begin{alertblock}{Discussion Question}%
}{%
  \end{alertblock}%
}
\newenvironment{discusion}{%
  \begin{alertblock}{Discussion Question}%
}{%
  \end{alertblock}%
}

\newenvironment{reflection}{%
  \begin{alertblock}{Classroom Reflection}%
}{%
  \end{alertblock}%
}
\newenvironment{reflexion}{%
  \begin{alertblock}{Classroom Reflection}%
}{%
  \end{alertblock}%
}

\newenvironment{paradox}{%
  \begin{alertblock}{Exploratory Paradox}%
}{%
  \end{alertblock}%
}
\newenvironment{paradoja}{%
  \begin{alertblock}{Exploratory Paradox}%
}{%
  \end{alertblock}%
}

\newenvironment{motivation}{%
  \begin{exampleblock}{Motivation: Data Science in Practice}%
}{%
  \end{exampleblock}%
}
\newenvironment{motivacion}{%
  \begin{exampleblock}{Motivation: Data Science in Practice}%
}{%
  \end{exampleblock}%
}

\newenvironment{quickchallenge}{%
  \begin{block}{Quick Team Challenge}%
}{%
  \end{block}%
}
\newenvironment{retoclase}{%
  \begin{block}{Quick Team Challenge}%
}{%
  \end{block}%
}


\title{PDF and Continuous Support}
\subtitle{Section 04.01 --- Density Function, Continuous Expectation and LOTUS}
\author[J. Castillo Colmenares]{Juliho Castillo Colmenares}
\institute[Tec de Monterrey]{Tecnologico de Monterrey}
\date{\vspace{-1.5cm}}

\begin{document}

% SLIDE 01: Title Page
\begin{frame}[plain]
	\titlepage
\end{frame}

% SLIDE 02: Roadmap
\begin{frame}{Roadmap --- Chapter 04: Continuous Random Variables}
	\vspace{-0.15cm}
	\begin{block}{Curricular Structure of Unit 3}
		\footnotesize
		\begin{enumerate}\itemsep=0.03cm
			\item \textbf{\textcolor{TecRojo}{Section 04.01: PDF and Continuous Support}} $\leftarrow$ \emph{Current focus}
			\item \textcolor{gray}{Section 04.02: Continuous CDF and Quantiles}
			\item \textcolor{gray}{Section 04.03: Expectation, Variance, and Continuous LOTUS}
			\item \textcolor{gray}{Section 04.04: Continuous Uniform Distribution}
			\item \textcolor{gray}{Section 04.05: Exponential Distribution and Memoryless Processes}
			\item \textcolor{gray}{Section 04.06: Normal Distribution and Z-Score}
			\item \textcolor{gray}{Section 04.07: Gamma, Beta, and Weibull Distributions}
		\end{enumerate}
	\end{block}
	\vspace{-0.2cm}
	\begin{alertblock}{Learning Objective}
		\scriptsize
		Introduce the Probability Density Function (PDF) as the continuous analogue of the PMF, master the normalization condition $\int f = 1$, compute probabilities via integration, and establish the foundations of continuous expectation and the LOTUS theorem.
	\end{alertblock}
\end{frame}

% SLIDE 03: Formal Definition and CDF
\begin{frame}{Formal Definition: PDF and Continuous CDF}
	\vspace{-0.15cm}
	\begin{columns}[T]
		\begin{column}{0.48\textwidth}
			\begin{block}{Probability Density Function (PDF)}
				\footnotesize
				A continuous random variable $X$ has PDF $f: \mathbb{R} \to [0, \infty)$ with $f(x) \ge 0$ and $\int f(x)\,dx = 1$. Probabilities are areas:
				$$P(a \le X \le b) = \int_{a}^{b} f(x)\,dx$$
			\end{block}
		\end{column}
		\begin{column}{0.48\textwidth}
			\begin{alertblock}{Continuous CDF}
				\footnotesize
				$F(x) = \int_{-\infty}^{x} f(t)\,dt$ with $F(-\infty) = 0$ and $F(+\infty) = 1$. By the Fundamental Theorem of Calculus, $f(x) = F'(x)$.
			\end{alertblock}
		\end{column}
	\end{columns}
\end{frame}

% SLIDE 04: Support and Continuous LOTUS
\begin{frame}{Continuous Support, Expectation, and LOTUS}
	\vspace{-0.15cm}
	\begin{columns}[T]
		\begin{column}{0.48\textwidth}
			\begin{block}{Continuous Support}
				\footnotesize
				$\mathcal{S}_X = \{x \in \mathbb{R} : f(x) > 0\}$ can be bounded $[a,b]$, semi-bounded $[a, \infty)$ or $(-\infty, b]$, or all $\mathbb{R}$. Examples: Uniform$([0,1])$, Exp$(\lambda) \to [0, \infty)$, Normal $\to \mathbb{R}$.
			\end{block}
		\end{column}
		\begin{column}{0.48\textwidth}
			\begin{alertblock}{Continuous LOTUS}
				\footnotesize
				$\E[g(X)] = \int g(x) f(x)\,dx$ under regularity conditions. Special cases: $g(x) = x$ gives $\E[X]$; $g(x) = (x - \mu)^{2}$ gives $\Var(X) = \E[X^{2}] - (\E[X])^{2}$.
			\end{alertblock}
		\end{column}
	\end{columns}
\end{frame}

% SLIDE 05: Applications
\begin{frame}{Applications in Engineering, Science, and Data Science}
	\vspace{-0.1cm}
	\begin{itemize}\itemsep=0.1cm
		\item \textbf{Telecommunications:} Packet waiting times, Gaussian channel noise, and interference models. \pause
		\item \textbf{Industrial Reliability:} Time-to-failure of mechanical and electronic components, exponential/Weibull modeling. \pause
		\item \textbf{Experimental Physics:} Measurement errors with Normal distribution, uncertainty propagation. \pause
		\item \textbf{Finance and Risk:} Brownian motion for asset prices, option valuation models.
	\end{itemize}
\end{frame}

% SLIDE 06: Python Lab Block 1
\begin{frame}[fragile]{Python Lab: PDF and Normalization}
	\vspace{-0.05cm}
	\scriptsize
	Computational verification of $\int f = 1$ for Exponential, Rayleigh, and Quadratic PDFs (\texttt{04.01\_pdf\_and\_support.py}):
	\vspace{0.02cm}
	{\lstset{basicstyle=\fontsize{5pt}{6pt}\selectfont\ttfamily}
	\lstinputlisting[firstline=21, lastline=44]{../../code/04_variables_aleatorias_continuas/04.01_pdf_and_support.py}}
\end{frame}

% SLIDE 07: Python Lab Block 2
\begin{frame}[fragile]{Python Lab: CDF and Interval Probabilities}
	\vspace{-0.05cm}
	\scriptsize
	Numerical CDF computation and $P(a \le X \le b)$ via `scipy.integrate.quad`:
	\vspace{0.02cm}
	{\lstset{basicstyle=\fontsize{4pt}{5pt}\selectfont\ttfamily}
	\lstinputlisting[firstline=47, lastline=73]{../../code/04_variables_aleatorias_continuas/04.01_pdf_and_support.py}}
\end{frame}

% SLIDE 08: Python Lab Block 3
\begin{frame}[fragile]{Python Lab: Continuous LOTUS and Moments}
	\vspace{-0.05cm}
	\scriptsize
	LOTUS verification: $\E[g(X)] = \int g(x) f(x)\,dx$ for Exponential and Rayleigh:
	\vspace{0.02cm}
	{\lstset{basicstyle=\fontsize{4.5pt}{5.5pt}\selectfont\ttfamily}
	\lstinputlisting[firstline=85, lastline=112]{../../code/04_variables_aleatorias_continuas/04.01_pdf_and_support.py}}
\end{frame}

% SLIDE 09: Python Lab Terminal Output
\begin{frame}[fragile]{Python Lab: Terminal Output}
	\vspace{-0.1cm}
	\begin{block}{}
		\fontsize{4pt}{5pt}\selectfont
		\begin{verbatim}
=== Block 1: PDF Validation & Normalization ===
  Exponential(0.5):         integral = 1.00000000
  Rayleigh(1.0):            integral = 1.00000000
  Quadratic(3/10):          integral = 1.00000000
  SciPy Rayleigh match: True

=== Block 2: CDF & Interval Probabilities ===
Exponential(0.5) CDF:
  F(0.5)=0.221199 | F(1.0)=0.393469 | F(2.0)=0.632121
P(0<=X<=1)=0.393469 | P(1<=X<=3)=0.383400
Rayleigh CDF: F(0.5)=0.117503 | F(1.0)=0.393469
Quadratic: F(-0.5)=0.2125 | F(0)=0.5 | F(0.5)=0.7875

=== Block 3: Continuous LOTUS & Moments ===
Exponential(0.5): E[X]=2.0 | E[X^2]=8.0 | Var=4.0
Rayleigh(1.0): E[X]=1.253314 | Var=0.429204
LOTUS E[sqrt(X)]: numerical=1.253314 | SciPy=1.253314
		\end{verbatim}
	\end{block}
\end{frame}

% SLIDE 10: Exercise Level 1 (Statement)
\begin{frame}{In-Class Exercise --- Level 1: Fundamental (1/2)}
	\vspace{-0.1cm}
	\begin{block}{Problem 4.1.1 --- Exponential PDF (Statement)}
		Consider $f(x) = c \cdot e^{-x}$ for $x \ge 0$ and $f(x) = 0$ for $x < 0$. Determine the value of $c$ such that $f$ is a valid PDF, calculate $P(1 \le X \le 3)$, and the CDF $F(x)$.
	\end{block}
	\vspace{0.15cm}
	\pause
	\begin{alertblock}{Model Setup}
		\begin{itemize}\itemsep=0.04cm
			\item Support: $[0, \infty)$ (semi-bounded).
			\item Normalization: $\int_{0}^{\infty} c \cdot e^{-x}\,dx = 1$.
			\item Probability: $P(1 \le X \le 3) = \int_{1}^{3} c \cdot e^{-x}\,dx$.
		\end{itemize}
	\end{alertblock}
\end{frame}

% SLIDE 11: Exercise Level 1 (Solution)
\begin{frame}{In-Class Exercise --- Level 1: Fundamental (2/2)}
	\vspace{-0.05cm}
	\scriptsize
	Evaluate the normalization integral and derive the CDF: \pause
	\begin{block}{Normalization Constant}
		\begin{align*}
			\int_{0}^{\infty} c \cdot e^{-x}\,dx = c \cdot [-e^{-x}]_{0}^{\infty} = c \cdot 1 = c = 1.
		\end{align*}
		Therefore $c = 1$ and the PDF is the standard exponential with $\lambda = 1$.
	\end{block}
	\vspace{0.05cm}
	\pause
	\begin{alertblock}{Probability and CDF}
		\scriptsize
		\begin{align*}
			P(1 \le X \le 3) &= \int_{1}^{3} e^{-x}\,dx = e^{-1} - e^{-3} \approx 0.3679 - 0.0498 = 0.3181, \\
			F(x) &= \int_{0}^{x} e^{-t}\,dt = 1 - e^{-x} \quad \text{for } x \ge 0.
		\end{align*}
		This is the CDF of the exponential distribution with $\lambda = 1$.
	\end{alertblock}
\end{frame}

% SLIDE 12: Exercise Level 2 (Statement)
\begin{frame}{In-Class Exercise --- Level 2: Operational (1/2)}
	\vspace{-0.1cm}
	\begin{block}{Problem 4.1.4 --- Rayleigh Distribution (Statement)}
		Let $X$ have CDF $F(x) = 1 - e^{-x^{2}/2}$ for $x \ge 0$. Derive the PDF, verify normalization, and calculate the median.
	\end{block}
	\vspace{0.15cm}
	\pause
	\begin{alertblock}{Model Setup}
		\begin{itemize}\itemsep=0.04cm
			\item The CDF is differentiable at $x > 0$ with derivative $f(x) = F'(x)$.
			\item Verification: $\int_{0}^{\infty} f(x)\,dx = 1$ via substitution $u = x^{2}/2$.
			\item Median: $F(m) = 0.5$ implies $m = \sqrt{2 \ln 2}$.
		\end{itemize}
	\end{alertblock}
\end{frame}

% SLIDE 13: Exercise Level 2 (Solution)
\begin{frame}{In-Class Exercise --- Level 2: Operational (2/2)}
	\vspace{-0.05cm}
	\scriptsize
	Derive and verify step by step: \pause
	\begin{block}{PDF Derivation}
		\scriptsize
		\begin{align*}
			f(x) = F'(x) = \dfrac{d}{dx}\left( 1 - e^{-x^{2}/2} \right) = x e^{-x^{2}/2}, \quad x \ge 0.
		\end{align*}
	\end{block}
	\vspace{0.05cm}
	\pause
	\begin{alertblock}{Normalization and Median}
		\scriptsize
		With $u = x^{2}/2$, $du = x\,dx$: $\int_{0}^{\infty} x e^{-x^{2}/2}\,dx = \int_{0}^{\infty} e^{-u}\,du = 1$. \quad \checkmark

		The median $m$ satisfies $1 - e^{-m^{2}/2} = 0.5$, i.e., $e^{-m^{2}/2} = 0.5$. Taking logarithm:
		\begin{align*}
			m = \sqrt{2 \ln 2} \approx 1.1774.
		\end{align*}
	\end{alertblock}
\end{frame}

% SLIDE 14: Exercise Level 3 (Statement)
\begin{frame}{In-Class Exercise --- Level 3: Analytical (1/2)}
	\vspace{-0.1cm}
	\begin{block}{Problem 4.1.7 --- Continuous Expectation Derivation (Statement)}
		Let $X$ be a continuous random variable with PDF $f(x)$. Prove formally that $\E[X] = \int x f(x)\,dx$ starting from $\E[X] = \int x\,dF(x)$ and integration by parts under regularity conditions.
	\end{block}
	\vspace{0.1cm}
	\pause
	\begin{alertblock}{Deduction Strategy}
		\scriptsize
		Use the representation $F(x) = \int_{-\infty}^{x} f(t)\,dt$ and the condition $\lim_{x \to \pm\infty} xF(x) = 0$ (required for convergence of the expectation).
	\end{alertblock}
\end{frame}

% SLIDE 15: Exercise Level 3 (Solution)
\begin{frame}{In-Class Exercise --- Level 3: Analytical (2/2)}
	\vspace{-0.05cm}
	\scriptsize
	Use the operational definition and direct substitution: \pause
	\begin{block}{Operational Substitution}
		\scriptsize
		By Leibniz rule and the definition $dF(x) = f(x)\,dx$:
		\begin{align*}
			\E[X] = \int_{-\infty}^{\infty} x \cdot dF(x) = \int_{-\infty}^{\infty} x \cdot f(x)\,dx.
		\end{align*}
		The technical condition $\int |x| f(x)\,dx < \infty$ guarantees the absolute convergence of the integral, equivalent to the light-tail condition $\lim_{|x| \to \infty} xF(x) = 0$. \quad \qedhere
	\end{block}
\end{frame}

% SLIDE 16: Exercise Level 4 (Statement)
\begin{frame}{In-Class Exercise --- Level 4: Challenging (1/2)}
	\vspace{-0.1cm}
	\begin{block}{Problem 4.1.10 --- Convolution of Gaussians (Statement)}
		Prove that if $X, Y \sim N(0, 1)$ are independent, then $Z = X + Y$ has PDF $N(0, 2)$. Generalize: if $X_{1}, \dots, X_{n}$ are i.i.d. $N(0, \sigma^{2})$, then $\sum X_{i} \sim N(0, n\sigma^{2})$.
	\end{block}
	\vspace{0.15cm}
	\pause
	\begin{alertblock}{Application}
		\scriptsize
		Interpret the result in the context of Gaussian noise filtering in telecommunications: if we average $n$ independent noisy samples, the variance of the average is $\sigma^{2}/n$.
	\end{alertblock}
\end{frame}

% SLIDE 17: Exercise Level 4 (Solution)
\begin{frame}{In-Class Exercise --- Level 4: Challenging (2/2)}
	\vspace{-0.05cm}
	\scriptsize
	Apply the convolution and complete the square: \pause
	\begin{block}{Convolution of Two Gaussians}
		\scriptsize
		\begin{align*}
			h(z) &= \int_{-\infty}^{\infty} \dfrac{1}{\sqrt{2\pi}} e^{-x^{2}/2} \cdot \dfrac{1}{\sqrt{2\pi}} e^{-(z-x)^{2}/2}\,dx \\
			&= \dfrac{1}{2\pi} e^{-z^{2}/4} \int_{-\infty}^{\infty} e^{-(x - z/2)^{2}}\,dx = \dfrac{1}{\sqrt{2\pi \cdot 2}} e^{-z^{2}/4}.
		\end{align*}
		This is the PDF of $N(0, 2)$. \quad \qedhere
	\end{block}
	\vspace{0.05cm}
	\pause
	\begin{alertblock}{Generalization and Interpretation}
		\scriptsize
		By induction, $\sum_{i=1}^{n} X_{i} \sim N(0, n\sigma^{2})$. If we average $n$ noisy samples, $\bar{X} = (1/n)\sum X_{i} \sim N(0, \sigma^{2}/n)$. The noise power is reduced by factor $n$, a fundamental principle of averaging filters in telecommunications.
	\end{alertblock}
\end{frame}

% SLIDE 18: Closing and Next Steps
\begin{frame}{Closing Section and Modular Perspective}
	\vspace{-0.2cm}
	\begin{block}{Synthesis: The Bridge from Discrete to Continuous}
		\scriptsize
		The Probability Density Function $f(x)$ is the continuous analogue of the PMF: instead of mass at points, it distributes probability over a continuum. The formal properties are analogous (non-negativity, integral = 1), but the calculation techniques require integration instead of summation.
	\end{block}
	\vspace{0.05cm}
	\pause
	\begin{alertblock}{Key Lessons and Next Step}
		\begin{itemize}\itemsep=0.05cm
			\item \textbf{Normalization:} $\int f = 1$ guarantees total mass; PDFs are relative, not absolute.
			\item \textbf{Support:} Determines where to focus computation; can be bounded, semi-bounded, or unbounded.
			\item \textbf{LOTUS:} $\E[g(X)] = \int g(x) f(x)\,dx$ avoids deriving the distribution of $g(X)$.
			\item \textbf{Next Section:} \textcolor{TecRojo}{Continuous CDF and Quantiles}.
		\end{itemize}
	\end{alertblock}
\end{frame}

\end{document}
